\documentclass[10pt,a4paper]{amsart}
\usepackage[margin=19mm]{geometry}
\usepackage{amsmath,amssymb,mathtools}
\usepackage[scr=rsfs,cal=euler]{mathalfa}
\usepackage{xfrac}
\usepackage[unicode,hidelinks,bookmarks=false]{hyperref}
\newcommand{\ie}{i.e.\ }
\newcommand{\cf}{cf.\ }
\newcommand{\resp}{respectively\ }
\newcommand{\Subsection}{\subsection}
\providecommand{\nobreakdash}{\nobreak\hbox{-}}
\setlength{\emergencystretch}{2em}
\multlinegap0pt
\usepackage{amssymb}
\usepackage{xcolor}
\usepackage{bm}
\usepackage[shortlabels]{enumitem}
\newtheorem{lemma}{Lemma}
\title{Compute Efficiency and Serial Runtime Tradeoffs \\for Stochastic Momentum Methods}
\author{Depen Morwani, Alexandru Meterez, Pranav Nair, Sham Kakade}
\date{Source: arXiv:2606.19179v1 (CC BY 4.0)}
\begin{document}
\maketitle

\noindent\emph{Source and licence.} Compute Efficiency and Serial Runtime Tradeoffs \\for Stochastic Momentum Methods. Version arXiv:2606.19179v1 (CC BY 4.0), distributed by arXiv under the Creative Commons Attribution 4.0 International licence, http://creativecommons.org/licenses/by/4.0/, as recorded in the arXiv licence metadata for this exact version and shown on the abstract page https://arxiv.org/abs/2606.19179v1. Full licence notice: \texttt{licenses/arxiv-2606.19179-CC-BY-4.0.txt}.\par\smallskip

\noindent\textit{Editorial scope.} Counted unit: the article's lemma lem:gap-etabeta (source0.tex lines 1717--1725). The article's complete original proof of it is delivered with it (source0.tex lines 1727--1830, one \texttt{proof} environment of 104 lines). No mathematics is written, repaired or re-derived: the statement and the proof are the article's own lines, in the article's own notation, and no retained line is substituted or re-flowed.  Retained uncounted as the source-local context the proof needs: lines 1473--1482.  Nothing was omitted from any retained span, so the package carries no omission record.  No retained line contains a citation, so the wrapper carries no bibliography and no bibliography entry.  No retained line contains a figure environment, an image-inclusion command or drawing code, so no image is delivered and the package declares no asset dependency.  Editorial formatting only, in addition: an AMS \texttt{amsart} preamble that restates the article's own declarations and macros used by the retained lines, namely the environments \texttt{lemma} on separate counters, together with the standard packages those lines need.  The retained environments print in this document as Lemma 1 in order of appearance, whereas the article prints the same items with the numbering of its own full document.  The complete native source of the article as distributed in the arXiv e-print for this exact version is bundled unchanged as \texttt{sources/203123/source0.tex}.
\begin{lemma}[Stochastic spectral radius dominates deterministic]
\label{lem:det_upper_bound}
For every $(\eta,\beta,B)$ we have
\[
\rho\big(\mathcal{T}_{H,\eta,\beta,B}\big)
\ \ge\
\rho\big(\mathcal{T}^{(\infty)}\big)
\;=\; \rho(A)^2.
\]
\end{lemma}

\begin{lemma}[Gap in terms of $\eta$ and $\beta$]
\label{lem:gap-etabeta}
For any stable HB--SGD parameters $(\eta,\beta,B)$, the spectral gap satisfies
\[
  1 - s(H,\eta,\beta,B)
  \;\le\;
  \min\big\{ 4 \eta \lambda_{\min} , 1-\beta \big\}.
\]
\end{lemma}

\begin{proof}
By Lemma~\ref{lem:det_upper_bound}, the stochastic covariance operator
$\mathcal{T}_{H,\eta,\beta,B}$ has spectral radius at least that of the
deterministic (full-batch) HB operator:
\[
s(H,\eta,\beta,B)
\;\ge\;
\rho\big(\mathcal{T}^{(\infty)}\big)
=
\rho(A)^2,
\]
where $A=\mathrm{blkdiag}(A_1,\dots,A_d)$ is the deterministic HB state
matrix in the $[w_t,w_{t-1}]$ state. Hence
\[
1 - s(H,\eta,\beta,B)
\ \le\
1 - \rho(A)^2.
\]

We first bound the gap by $1-\beta$. For each coordinate $i$, the
per–coordinate block $A_i$ has eigenvalues $r_{i,\pm}$ satisfying
$r_{i,+}r_{i,-}=\beta$. Thus
\[
\max\{|r_{i,+}|,|r_{i,-}|\}^2
\ \ge\ |r_{i,+}r_{i,-}| = \beta.
\]
Taking a maximum over $i$ yields $\rho(A)^2\ge\beta$, so
\[
1 - \rho(A)^2
\ \le\ 1 - \beta.
\]
Therefore
\[
1 - s(H,\eta,\beta,B)
\ \le\ 1 - \beta.
\]

We now show $1-s(H,\eta,\beta,B)\le 4\eta\lambda_{\min}$.
Let $A_{\min}$ be the $2\times2$ HB block corresponding to the smallest
eigenvalue $\lambda_{\min}$, with eigenvalues $r_{\min,\pm}$ solving
\[
z^2 - a_{\min} z + \beta = 0,
\qquad
a_{\min}=(1+\beta)-\eta(1-\beta)\lambda_{\min}.
\]
As $A_{\min}$ is a block of $A$, we have
\[
\rho(A)
\ \ge\ \max\{|r_{\min,+}|,|r_{\min,-}|\}.
\]
Hence
\[
1 - s(H,\eta,\beta,B)
\ \le\ 1 - \rho(A)^2
\ \le\ 1 - \max\{|r_{\min,+}|,|r_{\min,-}|\}^2.
\]

We distinguish the real and complex cases for $r_{\min,\pm}$.


For case of real roots (the overdamped case),
Assume $r_{\min,\pm}\in\mathbb{R}$, which corresponds to $a_{\min}^2>4\beta$.
We have $1 - \rho(A)^2 \leq (1-r_{\min,+}^2) \le 2(1-r_{\min,+})$.
The roots satisfy
\[
(r_{\min,+}-1)(r_{\min,-}-1)
= 1 - (r_{\min,+}+r_{\min,-}) + r_{\min,+}r_{\min,-}
= 1 - a_{\min} + \beta
= \eta(1-\beta)\lambda_{\min}.
\]
Since the roots are real and satisfy $r_+ r_- = \beta$ with $r_- \le
r_+$, we must have $r_- \le \sqrt{\beta}$. Using that $\beta\leq 1$ (stability),
\[
1 - \rho(A)^2 \leq 2(1-r_{\min,+}) = 2\frac{\eta(1-\beta)\lambda_{\min}}{1-r_-}
\le 2\frac{\eta(1-\beta)\lambda_{\min}}{1-\sqrt{\beta}}
= 2\eta(1+\sqrt{\beta})\lambda_{\min}
\leq 4\eta\lambda_{\min},
\]
which completes the proof of this case.

For case of complex roots (underdamped case),
assume $r_{\min,\pm}$ are complex conjugates, where $a_{\min}^2<4\beta$.
Then both have modulus
$\sqrt{\beta}$, so
\[
1 - \rho(A)^2
\ \le\ 1 - |r_{\min,+}|^2
= 1 - \beta.
\]
For the underdamped case (complex roots), due to that
$a_{\min}=(1+\beta)-\eta(1-\beta)\lambda_{\min}$, the condition
$a_{\min}^2<4\beta$ is equivalent to:
\[
  \eta(1-\beta)\lambda_{\min}  > (1-\sqrt{\beta})^2
\quad  \implies \quad
\eta(1+\sqrt{\beta})\lambda_{\min}  > 1-\sqrt{\beta}.
\]
Therefore,
\[
1-\beta \ =\ (1-\sqrt{\beta})(1+\sqrt{\beta})
\leq (1+\sqrt{\beta})^2\eta\lambda_{\min} \ \leq 4\eta\lambda_{\min},
\]
which completes the proof of the complex case.
\end{proof}

\end{document}
